% A template written for Examazej 1.7 (no app_version) with tags that 1.7 read otherwise
% than they are written: after an open-ended question, on the line of \item or %{#q}, two on
% one line, followed by a comment, with a small option letter.
%{#setting}
%university=Test University
%department=Mathematics
%coursecode=MATH 201
%examname=Quiz 2
%numberofvestions=4
%groups=3,(3),[3],2,1
%seed=old-layouts-seed
%{/setting}
\begin{enumerate}
\item
%{#q}Describe the graph of $y = x^2$.%{/q}
\item
%{#fixed}
%{#q}The vertex of $y = (x-1)^2$ is at%{/q}
\begin{enumerate}
\item %{#o}$(1, 0)$%{/o}
\item %{#o}$(0, 1)$%{/o}
\item %{#o}$(-1, 0)$%{/o}
\item %{#o}$(0, -1)$%{/o}
\end{enumerate}
\item %{#fixed}
%{#q}The axis of symmetry of $y = x^2 - 4x$ is%{/q}
\begin{enumerate}
\item %{#o}$x = 2$%{/o}
\item %{#o}$x = -2$%{/o}
\item %{#o}$x = 4$%{/o}
\end{enumerate}
\item
%{#fixed-options:c}
%{#q}Using the table above, $f(2)$ is%{/q}
\begin{enumerate}
\item %{#o}$1$%{/o}
\item %{#o}$2$%{/o}
\item %{#o}$4$%{/o}
\item %{#o}$8$%{/o}
\end{enumerate}
\item
%{#q}Explain how you read the table.%{/q}
\item
%{#fixed-options:B}
%{#q}Then $f(3)$ is%{/q}
\begin{enumerate}
\item %{#o}$6$%{/o}
\item %{#o}$16$%{/o}
\item %{#o}$9$%{/o}
\item %{#o}$12$%{/o}
\end{enumerate}
\item
%{#fixed} % keep this one in place
%{#q}The derivative of $x^3$ is%{/q}
\begin{enumerate}
\item %{#o}$3x^2$%{/o}
\item %{#o}$x^2$%{/o}
\item %{#o}$3x$%{/o}
\item %{#o}$x^3$%{/o}
\end{enumerate}
\item
%{#fixed} %{#q}The integral of $2x$ is%{/q}
\begin{enumerate}
\item %{#o}$x^2 + C$%{/o}
\item %{#o}$2x^2 + C$%{/o}
\item %{#o}$x + C$%{/o}
\end{enumerate}
\item
%{#fixed-options:A}
%{#q}$\lim_{x \to 0} \frac{\sin x}{x}$ equals%{/q}
\begin{enumerate}
\item %{#o}$1$%{/o}
\item %{#o}$0$%{/o}
\item %{#o}$\infty$%{/o}
\item %{#o}None of the above%{/o}
\end{enumerate}
\item
%{#q}$e^0$ equals%{/q}
\begin{enumerate}
\item %{#o}$1$%{/o}
\item %{#o}$0$%{/o}
\item %{#o}$e$%{/o}
\end{enumerate}
\item
%{#fixed}
%{#q}$\ln 1$ equals%{/q}
\begin{enumerate}
\item %{#o}$0$%{/o}
\item %{#o}$1$%{/o}
\item %{#o}$e$%{/o}
\end{enumerate}
\item
%{#separate-page} %{#fixed}
%{#q}Is every continuous function differentiable?%{/q}
\begin{enumerate}
\item %{#o}No%{/o}
\item %{#o}Yes%{/o}
\end{enumerate}
\end{enumerate}
